\chapter{Permutaciones, variaciones y combinaciones}
\label{md:combinatoria}

\section{Permutaciones como aplicaciones}

Una \textbf{permutación} de un conjunto finito $A$ es una biyección
$\sigma:A\to A$. Una notación de dos filas coloca los elementos
de $A$ arriba y sus imágenes debajo.
La composición se lee de derecha a izquierda:
$(\sigma_2\sigma_1)(a)=\sigma_2(\sigma_1(a))$.

Las permutaciones forman el \textbf{grupo simétrico} $S_A$:
la composición es asociativa, la identidad es el neutro y cada
permutación tiene inversa. En general, la composición no conmuta.
Para un conjunto con $n$ elementos se escribe también $S_n$.
Sus permutaciones se cuentan eligiendo sucesivamente las imágenes:
\[
 |S_n|=n(n-1)\cdots1=n!,\qquad 0!=1.
\]
Una ordenación de $n$ elementos distintos equivale a una permutación
de sus posiciones.

Las simetrías geométricas que conservan una figura permutan sus
vértices y forman un grupo por composición. En el triángulo equilátero,
la identidad, las dos rotaciones no triviales y las tres reflexiones
producen todas las permutaciones de sus tres vértices.

\section{Elegir con orden o sin él}

Sea un conjunto de $n\ge1$ elementos y un tamaño $r$.
Una \textbf{variación sin repetición} es una lista ordenada de $r$
elementos distintos: equivale a una aplicación inyectiva de
$\{1,\ldots,r\}$ al conjunto. Para $0\le r\le n$,
\[
 V(n,r)=n(n-1)\cdots(n-r+1)=\frac{n!}{(n-r)!}.
\]
Una \textbf{variación con repetición} permite volver a elegir un
elemento; equivale a una aplicación cualquiera. Para $r\ge0$,
\[
 VR(n,r)=n^r.
\]
En ambos casos, intercambiar posiciones puede cambiar el resultado.

Una \textbf{combinación sin repetición} es un subconjunto de $r$
elementos: solo importa qué elementos aparecen. Para $0\le r\le n$,
\[
 C(n,r)=\binom nr=\frac{n!}{r!(n-r)!}.
\]
Cada subconjunto admite $r!$ ordenaciones, luego $V(n,r)=r!C(n,r)$.
El tamaño cero tiene una sola elección, la vacía.

\begin{ejemplo}
Entre cinco personas, elegir dos representantes sin cargos da
$\binom52=10$ posibilidades. Asignar a dos personas distintas
los cargos de presidencia y secretaría da $V(5,2)=20$:
intercambiar los cargos produce otro resultado.
\end{ejemplo}

Una cadena binaria de longitud $n$ con exactamente $r$ unos queda
determinada al elegir sus posiciones: hay $\binom nr$.
Contando todas las cadenas por su número de unos,
\[
 \sum_{r=0}^n\binom nr=2^n.
\]

\section{Combinaciones con repetición}
\label{md:combinaciones-repeticion}

Una \textbf{combinación con repetición} es una selección sin orden
en la que cada elemento puede aparecer varias veces.
Dos selecciones son iguales si cada elemento tiene la misma
\emph{multiplicidad}; compartir los tipos de elementos no basta.
Si $x_i$ es la multiplicidad del elemento $i$, se cuenta
\[
 x_1+\cdots+x_n=r,\qquad x_i\in\N_0.
\]
Para $n\ge1$, $r\ge0$, el número es
\[
 \boxed{CR(n,r)=\binom{n+r-1}{r}.}
\]
La correspondencia de \textbf{estrellas y separadores} representa los
$r$ objetos iguales con estrellas y divide sus grupos mediante $n-1$
separadores. Elegir las posiciones de las estrellas entre $n+r-1$
lugares determina todas las multiplicidades, incluidas las nulas.
También cuenta distribuciones de $r$ objetos indistinguibles en
$n$ cajas distintas. Si los objetos fueran distinguibles, el modelo
sería una asignación con repetición y el total sería $n^r$.

\section{Ordenaciones circulares}

Para $1\le r\le n$, al colocar $r$ elementos distintos elegidos
entre $n$ alrededor de una circunferencia, se consideran iguales
las ordenaciones obtenidas por rotación, pero no por reflexión.
Cada ordenación circular corresponde a $r$ listas lineales:
\[
 \binom nr(r-1)!=\frac{V(n,r)}r.
\]
Si se usan los $n$ elementos, hay $(n-1)!$.
La división por $r$ utiliza que los elementos son distintos;
no se aplica automáticamente si se permiten repeticiones.
Las ordenaciones lineales con multiplicidades fijadas se cuentan
mediante los coeficientes multinómicos de la sección~\ref{md:multinomial}.
